\documentclass[10pt,a4paper]{amsart}
\usepackage[margin=19mm]{geometry}
\usepackage{amsmath,amssymb,mathtools}
\usepackage[scr=rsfs,cal=euler]{mathalfa}
\usepackage{xfrac}
\usepackage[unicode,hidelinks,bookmarks=false]{hyperref}
\newcommand{\ie}{i.e.\ }
\newcommand{\cf}{cf.\ }
\newcommand{\resp}{respectively\ }
\newcommand{\Subsection}{\subsection}
\providecommand{\nobreakdash}{\nobreak\hbox{-}}
\setlength{\emergencystretch}{2em}
\multlinegap0pt
\usepackage{bm}
\usepackage{bbm}
\usepackage{dsfont}
\usepackage{xspace}
\usepackage{cancel}
\usepackage{mathtools}
\usepackage{amsthm}
\makeatletter
\@ifundefined{prop}{\newtheorem{prop}{Prop}}{}
\@ifundefined{lem}{\newtheorem{lem}[prop]{Lemma}}{}
\makeatother
\newcommand{\bG}{\mathbb{G}}
\newcommand{\bT}{\mathbb{T}}
\newcommand{\cL}{\mathcal{L}}
\newcommand{\rd}{\mathrm{d}}
\title[Stable Degeneration, Non-degenerate Forms, and Kaledin's Con]{Stable Degeneration, Non-degenerate Forms, and Kaledin's Conjecture}
\author{Chenyang Xu, Ziquan Zhuang, Henri Guenancia}
\date{Source: arXiv:2606.02401v2 (CC BY 4.0)}
\begin{document}
\maketitle

\noindent\emph{Source and licence.} Stable Degeneration, Non-degenerate Forms, and Kaledin's Conjecture. Version arXiv:2606.02401v2 (CC BY 4.0), distributed under the Creative Commons Attribution 4.0 International licence, as recorded in the arXiv licence metadata for this exact version and shown on the abstract page https://arxiv.org/abs/2606.02401v2. Full rights notice: licenses/arxiv-2606.02401-CC-BY-4.0.txt.\par\smallskip

\noindent\textit{Editorial scope.} Counted unit: the Lem whose source label is \texttt{l-Sym=Hamilton} in the article (source0.tex lines 2381--2384, source label \texttt{l-Sym=Hamilton}), delivered together with the article's own complete original proof of it (source0.tex lines 2385--2432, one \texttt{proof} environment of 48 lines). No mathematics is written, repaired or re-derived: statement and proof are the article's own text line for line. Required context retained uncounted: none. Every definition, display and label of those lines is retained unchanged. Omitted from this package and not redistributed: 1--2380, 2433--2931. Nothing inside a retained excerpt is omitted or altered: every retained body is delivered exactly as the article's lines are written, so no omission receipt of any kind accompanies this package. Citations: the retained excerpts carry no citation command at all, so no bibliography entry of the article is transcribed and no reference is redistributed. Reference adaptation: none was needed and none was made; every reference inside the delivered unit resolves inside it, so no unresolved reference or citation is produced. Printed slips kept literal: no printed word, symbol, hypothesis or number of the counted statement or of its proof was altered. Layout and class: the article's own class is \texttt{amsart} with packages including amsmath, latexsym, amsfonts, amssymb, amsthm, geometry, hyperref, mathrsfs, graphicx, color, amsrefs, tikz, tikz-cd, bbm; the wrapper is the lane's pinned \texttt{amsart} editorial wrapper at 10pt on A4, which restates the theorem-like environments the retained text needs and adds the article's own macro definitions that the retained text needs (\texttt{\textbackslash bG}, \texttt{\textbackslash bT}, \texttt{\textbackslash cL}, \texttt{\textbackslash rd}), with extra wrapper packages (bm, bbm, dsfont, xspace, cancel, mathtools). The delivered numbering is the unit's own. Assets: no retained span contains a figure environment, a graphics-inclusion command, a PDF-page inclusion, a table or a TikZ picture; no image file is copied into this package, the package declares no asset dependency and no image file is redistributed with it. Licence: version arXiv:2606.02401v2 of this article is distributed under the Creative Commons Attribution 4.0 International licence, as recorded in the arXiv licence metadata for this exact version and shown on the abstract page https://arxiv.org/abs/2606.02401v2. A full rights notice is delivered at licenses/arxiv-2606.02401-CC-BY-4.0.txt. Characters: every retained excerpt is pure ASCII, so the delivered engine, XeLaTeX with its default Latin Modern text font, reports no missing character.
\begin{lem}\label{l-Sym=Hamilton}
% For a conical symplectic singularity $(Y,\sigma_Y)$, 
Any symplectic torus $\bT$-action on a conical symplectic singularity $(y\in Y,\sigma)$ that commutes with the conical $\bG_m$-action is Hamiltonian. 
\end{lem}

\begin{proof}
Assume the $\mathbb{G}_m$-action on $\sigma$ has weight $\ell>0$.
It is enough to construct the  Hamiltonian function for each
\(\xi\in\mathfrak t=\operatorname{Lie}(\mathbb T)\). Let $\xi_Y$ denote the corresponding vector field on $Y$. Since the $\mathbb T$-action is symplectic, it preserves $\sigma$ and we have $\mathcal L_{\xi_Y}\sigma=0$.
Hence, by Cartan's formula,
\[
\rd(\iota_{\xi_Y}\sigma)=\mathcal L_{\xi_Y}\sigma-\iota_{\xi_Y}\rd\sigma=0.
\]
Thus $\alpha_\xi:=\iota_{\xi_Y}\sigma$ is a closed regular \(1\)-form.

Let $\eta$ be the Euler vector field corresponding to the conical $\bG_m$-action. Then $\cL_{\eta}\sigma = \ell\sigma$. Because the $\mathbb T$-action commutes with the conical
\(\mathbb G_m\)-action, we have $ [\eta,\xi_Y]=0$. Therefore
\[
\mathcal L_\eta\alpha_\xi = \mathcal L_\eta(\iota_{\xi_Y}\sigma) = \iota_{[\eta,\xi_Y]}\sigma+\iota_{\xi_Y}(\mathcal L_\eta\sigma) = \ell\,\iota_{\xi_Y}\sigma =\ell\alpha_\xi .
\]
On the other hand, applying Cartan's formula to the closed \(1\)-form
\(\alpha_\xi\), we get
\[
\mathcal L_\eta\alpha_\xi
=
\rd(\iota_\eta \alpha_\xi)+\iota_\eta(\rd\alpha_\xi)
=
\rd(\iota_\eta\alpha_\xi).
\]
%Combining the two identities gives $\ell\alpha_\xi=d(\iota_E\alpha_\xi)$. 
Since $\ell>0$, combining the two identities gives %it follows that
\[
\alpha_\xi
=
\rd\left(\frac{1}{\ell}\iota_\eta\alpha_\xi\right)
=
\rd\left(\frac{1}{\ell}\iota_\eta\iota_{\xi_Y}\sigma\right).
\]
Thus the Hamiltonian function for \(\xi\) is
\[
H_\xi
=
\frac{1}{\ell}\iota_\eta\iota_{\xi_Y}\sigma
=
\frac{1}{\ell}\sigma(\xi_Y,\eta),
\]
up to the usual sign convention for contractions. % Since the $\bT$-action preserves $\sigma$ and commutes with both $\xi_Y$ and $\eta$, the function $H_\xi$ is $\bT$-invariant. 
In particular,
\[
\rd H_\xi=\iota_{\xi_Y}\sigma \, ,
\]
so the action is Hamiltonian.
\end{proof}

\end{document}
